\chapter{Industrialisation in India: Caste, Plantation, Mining and the Informal Sector}

\section{Caste in the Mills: Morris and Chandavarkar}

\begin{KeyConcept}
\begin{itemize}
\item \textbf{Morris D. Morris} (Bombay textile mills): recruitment was \textbf{not} based on caste, but \textbf{Dalits were excluded from the spinning section} --- by \textbf{both Hindus and Muslims} --- because spinning involved putting a thread in the needle using one's \textbf{spit (saliva)}, and there was the idea of pollution. The spinning section was the \textbf{highest-paid}, so it was monopolised for one's own caste and kin. Dalit migration to industry was \textbf{late} (they initially stayed in the villages, then came for cleaning/manual work).
\item \textbf{Rajnarayan Chandavarkar} (industries 1900--1940) complemented Morris by studying \textbf{settlement patterns}: workers of different castes lived in \textbf{clusters} in common buildings (the one-room \textbf{chawls}, in semi-urbanised areas), with \textbf{very little interaction} between castes outside the workplace. Alongside the large factories, \textbf{small businesses (petty bourgeois)} mushroomed --- paying low wages and outside state regulation.
\item The workers \textbf{compartmentalised} their lives: at the workplace there is \textbf{organic solidarity} (specialised division of labour, cooperation across castes), but outside there is \textbf{no mechanical solidarity} (they return to caste, with little interaction). This is \textbf{Milton Singer's} critique of Weber's Protestant ethic --- the \textbf{Chettiars} of South India (among Asia's first bankers, with the \textbf{hundi} system) compartmentalise personal and professional life, so the spirit of capitalism is \textbf{not unique to Protestants}.
\end{itemize}
\end{KeyConcept}

\begin{QuickRevision}
\begin{itemize}
\item Morris: no caste-based recruitment, but Dalits excluded from spinning (spit/pollution; highest wage); late Dalit migration.
\item Chandavarkar: caste clusters in chawls; petty-bourgeois small businesses; compartmentalisation.
\item Organic solidarity at workplace, not mechanical outside; Milton Singer (Chettiars, hundi) vs Weber's Protestant ethic.
\end{itemize}
\end{QuickRevision}

\section{Plantation Labour: Indentured and Bonded}

\begin{KeyConcept}
\begin{itemize}
\item \textbf{Sarit Bhowmik} (a leading name in the sociology of work) studied the \textbf{tea plantations of Assam and West Bengal (Darjeeling)}. Assam is the oldest tea plantation. Workers were recruited largely from the \textbf{tribal population of Central India}.
\item \textbf{Indentured labour} (instituted \textbf{after the abolition of slavery}): a bonded-labour contract of \textbf{five years or more} (the labourers were called \textbf{Girmitiyas}, from 'girmit' = the agreement/permit). They were recruited for \textbf{sugar, cotton and tea plantations and rail construction} in the British colonies (West Indies, Africa, South-East Asia), derogatorily called \textbf{'coolies'}. Mostly illiterate, they \textbf{could not understand the contract} (they put thumb imprints). In \textbf{1856--57 the death rate en route was 10--13\%}; conditions were harsh (3 a.m. starts, physical torture, children working from age 5, desertion punished by extending the contract).
\item \textbf{Radhika Singha}: the labourers accepted these contracts to \textbf{escape poverty and famines} that were frequent under colonial rule.
\item In \textbf{West Bengal} there was no indentured labour, but \textbf{bondage through indebtedness} --- money-lenders (in collusion with the plantation owners) lent to the workers, who could then never leave. Wages were \textbf{about half} those of textile workers, with working hours about \textbf{1.5 times} longer.
\item About \textbf{20--25\% of plantation workers were women} and \textbf{5--10\% children}.
\end{itemize}
\end{KeyConcept}

\section{Mining: Janaki Nair and the Kolar Gold Fields}

\begin{KeyConcept}
\begin{itemize}
\item \textbf{Janaki Nair} studied the \textbf{Kolar Gold Fields} (the only gold mines in India) in colonial times. Most workers were \textbf{Dalits}, who had moved from neighbouring areas \textbf{as an escape from caste oppression} (they hoped caste would not operate in mining). But they found \textbf{bondage through debt} --- money-lenders, in collusion with the mine owners, lent to the workers and trapped them.
\item The workers \textbf{pressured the trade unions} to raise the issue of \textbf{untouchability} (along with wages, working hours and training) --- and they succeeded.
\item \textbf{Overall pattern}: different industries had different caste compositions --- \textbf{textile} = caste groups of the varna system (with Dalits excluded); \textbf{plantations} = tribals; \textbf{mining} = Dalits.
\item \textbf{Trade unions} consist only of workers (not executives or diploma-holders) --- a Marxist idea, since the interests of workers and managers are contradictory.
\end{itemize}
\end{KeyConcept}

\begin{QuickRevision}
\begin{itemize}
\item Plantation (Sarit Bhowmik): Assam/West Bengal tea; tribal workers; indentured labour (Girmitiyas, 5+ years, after abolition of slavery, 10--13\% death rate); Radhika Singha (escape poverty/famine); West Bengal bonded via debt; half wages, 1.5× hours; 20--25\% women.
\item Mining (Janaki Nair, Kolar): Dalit workers as escape from caste; debt bondage; trade unions raised untouchability.
\item Caste pattern: textile = varna castes, plantation = tribals, mining = Dalits.
\end{itemize}
\end{QuickRevision}

\section{Post-Colonial vs Post-Independence}

\begin{KeyConcept}
\textbf{Post-independence} simply means after 1947. \textbf{Post-colonial} is a \textbf{political term} (used by the \textbf{subaltern school}): though India won independence from British rule, \textbf{many pockets --- Dalits, women, the disabled, tribals --- are still colonised}, still exploited and oppressed, and often without even the sense of nationalism. So these scholars prefer 'post-colonial' to 'independent India'.
\end{KeyConcept}

\begin{QuickRevision}
Post-independence = after 1947; post-colonial = political term (subaltern school) --- many groups still colonised.
\end{QuickRevision}

\section{Labour Commitment: Moore, Feldman, Lambert and Sheth}

\begin{KeyConcept}
\begin{itemize}
\item \textbf{Wilbert Moore and Arnold Feldman} argued that India \textbf{could not industrialise} because its social structure impedes \textbf{labour commitment}: a \textbf{closed system of stratification (caste)}, emphasis on \textbf{primordial loyalties} (caste, region, religion, language), the high role of \textbf{religious values}, and a \textbf{strong attachment to land} (kept as a status symbol and economic back-up).
\item \textbf{Richard D. Lambert} (five factories in Pune) found the \textbf{opposite}: workers in small factories looked for big-factory jobs, and once in a big factory showed \textbf{lifetime commitment} --- they were \textbf{over-committed}, viewing the job the way they viewed their \textbf{traditional caste occupation} (the specialist serves the patron --- the jajmani system).
\item \textbf{N. R. Sheth} (a factory, early 1960s) found that \textbf{traditional culture can promote commitment}: recruitment was based on \textbf{particularistic norms of obligation to caste and kin}, and workers accepted their obligation to supervisors \textbf{as religious duty} --- so there is \textbf{no contradiction between traditional values and industrialism}.
\end{itemize}
\end{KeyConcept}

\begin{QuickRevision}
\begin{itemize}
\item Moore \& Feldman: India can't industrialise (closed stratification, primordial loyalties, religion, attachment to land).
\item Lambert (Pune): Indian workers over-committed (lifetime, job = caste occupation/jajmani).
\item Sheth: traditional culture promotes commitment (particularistic recruitment; obligation as religious duty).
\end{itemize}
\end{QuickRevision}

\section{The Informal Economy: Keith Hart, Breman and Sanyal}

\begin{KeyConcept}
\begin{itemize}
\item \textbf{Keith Hart} coined the concept of the \textbf{informal economy} --- the state of economic activities of the \textbf{working poor}, with an \textbf{absence of social security provisions}, based on \textbf{rural--urban migration}.
\item \textbf{Jan Breman's} characteristics of the informal sector: \textbf{capital comes from kinship groups}; \textbf{labour comes from kinship + caste groups} (and the village/neighbourhood --- the owner personally knows the workers, who work out of \textbf{personal obligation}); \textbf{skills are learned on the workplace} (no prior training); and there is \textbf{deviance in production} --- pirated goods (Gaffar Market, Tank Road, Karol Bagh --- the largest piracy markets, assembling pirated electronics in Delhi's slums).
\item \textbf{Ravi Sundaram} ('Pirate Modernity', CSDS): around the Asian Games, when TVs became popular, people went to China and returned with equipment to \textbf{make pirated televisions in Delhi's slums} --- a 'pirate modernity'.
\item \textbf{Kalyan Sanyal}: the informal sector is the \textbf{space where people innovate new forms of production, distribution and redistribution} to accommodate the \textbf{surplus labour force} at levels of economic subsistence.
\item \textbf{Over 90\% of India's labour force works in the informal sector.}
\end{itemize}
\end{KeyConcept}

\section{Post-LPG Informalisation and Class Structure}

\begin{KeyConcept}
\begin{itemize}
\item After \textbf{LPG liberalisation}, economists expected \textbf{formalisation}, but the opposite happened --- \textbf{informalisation} of the labour force. \textbf{Kalyan Sanyal and Rajesh Bhattacharya} give two causes: (1) the \textbf{dismantling of trade and capital barriers} --- MNCs came with better technology, so Indian companies \textbf{outsourced to the informal sector} (cheaper) and small/medium enterprises shut down (about \textbf{1 lakh textile workers} were removed after liberalisation, of whom \textbf{80,000 joined the informal economy}); (2) \textbf{cheaper imports} (the toy industry was finished by Chinese imports). Critics thus say globalisation did not benefit the mass of the population.
\item \textbf{Class structure of the informal sector}: (1) \textbf{owners} --- a \textbf{petty bourgeois} class (mini-workshops, self-employed artisans); (2) \textbf{sub-proletariat} --- the \textbf{largest} segment, casual + unskilled workers bound by kinship ties and personal obligation; (3) the \textbf{declassed} --- those who have severed ties with family and village, live like vagabonds, work irregularly, and take up deviant activities (selling pirated goods), lacking fixed accommodation and the stamina for full-time work (due to drugs).
\item \textbf{Class structure of the formal economy}: (1) the \textbf{capitalist class} (control over the means of production); (2) the \textbf{professional class} (from the 1960s, trained in technical institutes --- white-collar, high wage, high status, social security, not from a working-class background, mostly upper caste --- first entering PSUs and engineering); (3) the \textbf{proletariat} (from the rural middle class, in trade unions, owning land in the village, indifferent or opposed to agricultural labourers); (4) \textbf{manual workers} (from the rural lower class, low in the caste hierarchy).
\end{itemize}
\end{KeyConcept}

\begin{QuickRevision}
\begin{itemize}
\item Informal economy (Keith Hart): working poor, no social security, rural--urban migration.
\item Jan Breman: capital/labour from kinship; personal obligation; on-the-job skill; deviance (piracy --- Gaffar Market; Ravi Sundaram's 'Pirate Modernity').
\item Kalyan Sanyal: informal = innovation to absorb surplus labour; 90\%+ of labour force.
\item Post-LPG informalisation (Sanyal \& Bhattacharya): MNC competition $\rightarrow$ outsourcing; cheaper imports (toys); 1 lakh textile workers, 80,000 to informal.
\item Informal class structure: petty-bourgeois owners; sub-proletariat (largest); declassed.
\item Formal class structure: capitalist; professional (1960s, upper caste); proletariat (rural middle class, trade unions); manual workers (rural lower class).
\end{itemize}
\end{QuickRevision}

